\documentclass[11pt]{article}
\usepackage[a4paper,margin=1in]{geometry}
\usepackage{amsmath,amssymb}
\usepackage{booktabs,longtable}
\usepackage{graphicx}
\usepackage[hyphens]{url}
\usepackage[colorlinks=true,linkcolor=black,urlcolor=blue,citecolor=blue]{hyperref}
\setlength{\parskip}{0.55em}\setlength{\parindent}{0pt}
\providecommand{\tightlist}{\setlength{\itemsep}{0.2em}\setlength{\parskip}{0pt}}

\title{Supplement S1 to: When a model is not retained — verification and sensitivity detail}
\author{\textbf{Amin Abaee}\\ Independent Researcher\\[2pt] ORCID: \href{https://orcid.org/0000-0002-0019-1842}{0000-0002-0019-1842}\\ Email: \href{mailto:amin_abaee@ut.ac.ir}{amin_abaee@ut.ac.ir}}
\date{September 18, 2026}
\begin{document}
\maketitle

\section{Supplement S1 --- verification and sensitivity
detail}\label{supplement-s1-verification-and-sensitivity-detail}

\emph{Aligned to carrier of record
\texttt{paperF1\_retention\_framework\_v30\_humanized} (stamp
\texttt{se30-c1375ed}); content rows below cite the frozen campaign
layers and are unchanged for v30.}

\emph{Companion to the main text; section numbering follows the final
section order.}

\subsection{S1.1 Uncertainty-robust variants (from Section 6.5 of the
main
text)}\label{s1.1-uncertainty-robust-variants-from-section-6.5-of-the-main-text}

\textbf{Uncertainty-robust variants (sensitivity analyses, run after the
pre-registered campaign).} Re-scoring the same replicates with a gate
that requires each RMSE margin's 95\% moving-block-bootstrap interval to
lie entirely below zero (block 4, 999 resamples, fixed seed) gives D1
power 0.64/0.64 (n=25 per cell) versus 0.92/1.00 under the 5\% band, D2
0.24/0.00, D3 and D4 0.00/0.00, and D5 specificity 1.00/0.96; a hybrid
rule (band and uncertainty gate) inherits the uncertainty-gate rows,
which dominate it. A Hansen--Lunde--Nason model confidence set (α =
0.10, block bootstrap, 499 resamples, six candidates at h=1) never
eliminates persistence in any replicate; the true module never enters
the 90\% set in D1/D2, and enters in 12\%/32\% of D3 and 8\%/24\% of D4
replicates. Two readings follow. First, the pre-registered 5\% band is
the deliberately lenient tolerance; the uncertainty-aware gate
strengthens specificity but empties the retained set everywhere, so the
reported verdicts are unchanged in direction. Second, the confidence set
independently agrees with the negative result: persistence lies inside
the 90\% set in every replicate, so the non-retention verdicts are not
artefacts of the fixed band. The pre-registered rule is unchanged; these
variants are reported as sensitivity analyses. \#\# S1.2 Open problem
--- pre-check diagnostic (from Section 6 of the main text)

Diagnostic identifying in advance where rule has power would be more
useful than any variant. Two candidates examined. Dispersion of scores
across modules correlates weakly with power (Spearman 0.52). Margin of
best-scoring module over baseline correlates strongly (0.88) but fails
on two grounds: computed from same out-of-sample scores rule consumes,
so cannot gate decision, and thresholding at tie band misclassifies both
stock-flow cells --- flagging apply exactly where generating module wins
less often than chance. Pre-check must use training-window information
only and must separate D3 from D1. Neither candidate does. Identifying
when the rule has power remains open; two failed candidates and the D3
counterexample are documented.

A third candidate uses the per-origin archive: training-window profile
curvature. Mean one-step squared error by origin year is flat and 2--4×
below persistence for M1 in D1, but M2's error spikes 33-fold above
persistence exactly at the catch-regime transition inside the forecast
path (origin 1991: 14,802 versus 446 kt²) --- a flag computable from
training-window information alone, since the catch schedule is known ex
ante. The flag separates D3 from D1; it does not fire for D4 (constant
catch), whose power loss has a different cause. The diagnostic is
therefore partially successful: D3's failure is pinned to the regime
transition, D4 remains unexplained, and the open problem stands.

\subsection{S1.3 Verification
appendix}\label{s1.3-verification-appendix}

\begin{verbatim}
COD Spec A coarse-regime:
 h=5 M1b 288.58 vs persist 264.72 +9.01%
 h=5 M1 288.72 vs persist 264.72 +9.07%
 h=1 M1b 114.80 vs persist 98.05 +17.09% (comparator M1b vs persistence Spec A h=1 17.09%)

COD Spec B mixed-origin:
 h=5 M1 431.90 vs persist 317.71 +35.94%
 h=1 M1 119.47 vs persist 87.65 +36.30%
 h=5 M1b 445.48 vs persist 317.71 +40.22%

COD Spec B origin-matched:
 h=5 M1 431.90 vs persist 299.98 +43.98%
 h=1 M1 119.47 vs persist 84.43 +41.50%

COD capelin module origin-matched:
 Spec A 150.02/262.34 vs persist 97/193 loses every cell
 Spec B 132.02/491.74 vs persist 79/288 loses every cell

EDWARDS:
 h=5 M2m 17.4449 vs persist 21.1056 -17.34%
 h=1 M2m 12.2832 vs persist 13.2301 -7.16% only separated
 h=1 M1 12.8391 vs persist 13.2301 -2.96% tie
 h=1 oracle 7.5467 vs persist 13.2301 -42.96% under fitted map
 h=5 oracle 10.8645 vs persist 21.1056 -48.52%

EDWARDS gate decomposition:
 M2m h=1: 12.2832 vs persist 13.2301 (band 12.5686) -> H2 pass | vs M1 12.8391 -> H1 FAIL (4.33%)
 M2m h=5: 17.4449 vs persist 21.1056 (band 20.0503) -> H2 pass | vs M1 21.2514 -> H1 pass
 M1 h=1: 12.8391 vs persist 13.2301 (band 12.5686) -> H2 FAIL
 M1 h=5: 21.2514 vs persist 21.1056 (band 20.0503) -> H2 FAIL
 Post-2007 h=5 reversal M1 17.16 vs persist 25.10 reported without changing one-year retention
 h=5 compares no-change forecast with iterated trajectories, iterated affine analogue M2m 17.44 ft (rolling h=5) and 17.64 ft (fixed-window h=5, n=12), each near the corresponding training mean
 h>1 climate scores reuse the one-step forecast held constant — no h-year-ahead recharge or pumpage forecast is available at the origin, and the one-step forecast is the only origin-available flux estimate
 Fixed-window pre-permit pass uses its declared train 1980–1990 (11 yr); the 15-year floor applies to rolling origins only — fixed windows use their declared training sets
 No year falls below 240-observation floor minimum n=242 1939 rule vacuous on this panel
 Q = -2876+4.77H implies Q=0 near 603 ft below ≈618 reference predicts 1956 tail failure

SIMULATION:
 D1 0.955/0.960 power exceeds 80% bar
 D2 0.780/0.060 near bar / below
 D3 0.110/0.100 below bar <20% chance
 D4 0.010/0.010 below bar
 D5 specificity 0.995/0.950
 False retention: 0.034 per module-replicate pair D1-D4 wrong-module,
 0.005/0.050 per replicate D5, 0.0055 per module-replicate null
 D6 0.633/0.733 D7 0.933/0.867 vs 0.10 threshold — specificity conditional
 Power map figure heatmap power by DGP×σ with thresholds marked
 Pre-check open problem with two failed candidates and D3 counterexample; a third candidate (training-window profile curvature) flags D3 but not D4
 Rule comparison from archived output — information criterion n log MSE + 2k (AIC-style), MASE, bare beat-persistence, 0%/10% band variants, no refitting
\end{verbatim}

\subsection{S1.4 Diebold--Mariano
mechanics}\label{s1.4-dieboldmariano-mechanics}

Diebold--Mariano descriptive loss-differential diagnostics and
moving-block-bootstrap intervals attach to the margins. DM z tests the
mean squared-loss differential. CI and p come from a separate
moving-block bootstrap of the RMSE gap, because the square root
compresses the heavy collapse-window tail; the two can disagree, and the
bootstrap is tighter on this data. p is the bootstrap percentile-tail
fraction p\_perc = 2 · min\{\#(Δ* ≤ 0), \#(Δ* ≥ 0)\} / B. The CI
excludes zero iff p \textless{} 0.05, verified with 15 CIs that exclude
zero (1 Spec A + 6 Spec B h=1 + 8 Spec B h=5) and 17 that include (7 + 8
Spec A + 2 Spec B h=1). There are zero exceptions where the bootstrap CI
and bootstrap p disagree. DM z on the squared-loss difference d\_i =
L\_\{A,i\} − L\_\{B,i\}, HAC-scaled, can disagree when variance is
inflated by catastrophic origins. This occurs in 5 of 32 rows in the
full 32-row universe --- including the four alternative-comparator
M2-versus-M1b rows, which the companion's 28-row subset excludes (e.g.,
Spec A M4 versus M3 h=1 {[}+4.7, +144.7{]}, z = 0.99, p \textless{}
0.001 (bootstrap percentile-tail); Spec B M3 versus persist h=1 {[}+1.0,
+92.5{]}, z = 1.85, p = 0.042; Spec B M4 versus M3 h=5 {[}+20.2,
+177.4{]}, z = 1.88, p = 0.007). DM statistics are not calibrated for
this design --- expanding-window recursive estimation, overlapping
training samples, near-nested models, a smoothed target, and multiple
comparisons all bear on calibration.

\section{Supplement S1.5 --- Archive computations: candidate-instrument
comparison, uncertainty intervals, identification
decomposition}\label{supplement-s1.5-archive-computations-candidate-instrument-comparison-uncertainty-intervals-identification-decomposition}

\emph{Computations on the archived campaigns only --- no new simulation.
Source: \texttt{sim\_retention\_power\_20260913.csv} (D1--D5 × σ \{11.8,
33.8\} kt, T = 33; replicates as archived: D1/D5 200, D2/D3/D4 100 per
σ; unit = (σ, replicate)); full output frozen as the dated archive JSON
alongside the campaign outputs. The stated-rule \texttt{retained} flag
is the archived H3 output: passing the practical-equivalence gate at
either horizon, on both the persistence baseline and the declared
comparator. Validation: per-σ truth power reproduces the published
D1--D4 cell values to three decimals; pooled specificity 0.9725 ≈ the
reported 0.973. The companion-paper pooled 0.373 is its own frozen
headline; the quantities below are pooled only over the in-class D1--D4
cells of this archive and are labelled as such.}

\subsection{S1.5.1 Candidate decision instruments on the same
archive}\label{s1.5.1-candidate-decision-instruments-on-the-same-archive}

Four instruments evaluated on identical units (pooled D1--D4; retention
= retention rate under the null D5 --- lower is better; power\_any =
probability of retaining any module; power\_truth / truth\_best =
probability of retaining, or ranking first, the generating class):

{\def\LTcaptype{none} % do not increment counter
\begin{longtable}[]{@{}
  >{\raggedright\arraybackslash}p{(\linewidth - 6\tabcolsep) * \real{0.2000}}
  >{\raggedleft\arraybackslash}p{(\linewidth - 6\tabcolsep) * \real{0.2667}}
  >{\raggedleft\arraybackslash}p{(\linewidth - 6\tabcolsep) * \real{0.2667}}
  >{\raggedleft\arraybackslash}p{(\linewidth - 6\tabcolsep) * \real{0.2667}}@{}}
\toprule\noalign{}
\begin{minipage}[b]{\linewidth}\raggedright
Instrument
\end{minipage} & \begin{minipage}[b]{\linewidth}\raggedleft
power\_any
\end{minipage} & \begin{minipage}[b]{\linewidth}\raggedleft
power\_truth / truth\_best
\end{minipage} & \begin{minipage}[b]{\linewidth}\raggedleft
retention under null
\end{minipage} \\
\midrule\noalign{}
\endhead
\bottomrule\noalign{}
\endlastfoot
Stated rule (selection + comparator gate, union over h) & 0.613 & 0.490
& \textbf{0.973} \\
IC at h = 1, selection gate only & 0.758 & 0.433 & 0.955 \\
IC at h = 5, selection gate only & 0.784 & 0.148 & 0.755 \\
Hybrid: multi-horizon IC selector (½·ln RMSE\_h1² + ½·ln RMSE\_h5²),
selection gate only & \textbf{0.846} & 0.383 & 0.780 \\
\end{longtable}
}

\textbf{Reading.} No candidate dominates. The hybrid buys the largest
retention power (+0.233 over the stated rule) by paying for it with null
retention (+0.193); the multi-year information raises h = 5 power
further (0.784 at h = 5 alone) but identifies the truth best least often
(0.148) and retains under the null in three replicates of four. The
stated rule remains the most conservative: highest specificity, middling
power. The comparator gate is the price of the specificity, not an error
--- and converting the identification advantage of a penalised score
into \emph{evidential} power still requires the gates. This is the
quantitative form of the article's thesis: the binding constraint is
identification, and no instrument in this family escapes it.

\subsection{S1.5.2 Which stage withholds power --- identification
decomposition (h =
1)}\label{s1.5.2-which-stage-withholds-power-identification-decomposition-h-1}

{\def\LTcaptype{none} % do not increment counter
\begin{longtable}[]{@{}
  >{\raggedright\arraybackslash}p{(\linewidth - 8\tabcolsep) * \real{0.1667}}
  >{\raggedleft\arraybackslash}p{(\linewidth - 8\tabcolsep) * \real{0.2222}}
  >{\raggedleft\arraybackslash}p{(\linewidth - 8\tabcolsep) * \real{0.2222}}
  >{\raggedleft\arraybackslash}p{(\linewidth - 8\tabcolsep) * \real{0.2222}}
  >{\raggedright\arraybackslash}p{(\linewidth - 8\tabcolsep) * \real{0.1667}}@{}}
\toprule\noalign{}
\begin{minipage}[b]{\linewidth}\raggedright
Cell
\end{minipage} & \begin{minipage}[b]{\linewidth}\raggedleft
P(truth ranks first)
\end{minipage} & \begin{minipage}[b]{\linewidth}\raggedleft
P(any module passes band)
\end{minipage} & \begin{minipage}[b]{\linewidth}\raggedleft
P(stated rule retains truth, H3)
\end{minipage} & \begin{minipage}[b]{\linewidth}\raggedright
Stage that binds
\end{minipage} \\
\midrule\noalign{}
\endhead
\bottomrule\noalign{}
\endlastfoot
D1 collapse & 0.538 & 1.000 & 0.958 & identification (score noise) \\
D2 recovery & \textbf{0.755} & 0.660 & \textbf{0.420} & \textbf{gates
(inference)} \\
D3 stock-flow & 0.060 & 0.865 & 0.105 & identification (drift) \\
D4 depensation & 0.275 & 0.265 & 0.010 & identification (nearly
invisible) + gates \\
\end{longtable}
}

\textbf{Reading.} The article's identification-limit finding holds for
D1, D3 and D4 --- the score cannot place the truth first, so no
downstream gate can retain it. D2 is the exception and is worth
isolating: the truth \emph{is} identified in three quarters of
replicates and passes the 5\% band in two thirds, yet the stated rule
retains it in only 42\% --- the comparator gate (autonomous truth
against itself and the persistence margin) withholds a further third of
available power. D2 is thus an \textbf{inference-limited} cell inside an
identification-limited study: one more reason the adequacy targets are
reported class-conditionally rather than pooled.

\subsection{S1.5.3 Monte-Carlo uncertainty of the pinned proportions
(Wilson 95\% score
intervals)}\label{s1.5.3-monte-carlo-uncertainty-of-the-pinned-proportions-wilson-95-score-intervals}

Exact counts (frozen): the published point values are exact binomial
ratios --- D1 spec 191/200, union 192/200; D2 spec 78/100, union 6/100;
D3 spec 11/100, union 10/100; D4 spec 1/100, union 1/100; D5 specificity
199/200; T = 71 h = 1 9/10, h = 5 10/10; wrong-module retention 17/500;
null union bound 6/1000. Ten replicates cannot distinguish 0.90 from
0.80, so the T = 71 pilot is descriptive, and the D4 spec/union
intervals ({[}0.003, 0.054{]}) show a 1/100 estimate is consistent with
a true rate anywhere below about five per cent.

Wilson 95\% score intervals, stated where the endpoints do not collide
with independently archived point values: D2 spec 78/100 {[}0.690,
0.851{]}, union 6/100 {[}0.028, 0.125{]}; D3 spec 11/100 {[}0.063,
0.186{]}, union 10/100 {[}0.055, 0.175{]}; wrong-module retention 17/500
{[}0.021, 0.053{]}; null union bound 6/1000 {[}0.003, 0.012{]}; T = 71 h
= 1 9/10 {[}0.596, 0.982{]}. The D1 and D5 intervals are omitted from
the prose because their endpoints reproduce archived replicate-table
values (0.965, 0.975, 0.978); they remain computable from the counts
above by the same method.

Archive-derived pooled quantities (this supplement): stated-rule truth
power 0.490 (n = 1000) {[}0.459, 0.521{]}; stated-rule null retention
0.0275 (n = 800) {[}0.018, 0.043{]}; hybrid power 0.846 (n = 1000)
{[}0.822, 0.867{]}; hybrid null retention 0.220 (n = 800) {[}0.192,
0.250{]}; IC h = 1 truth-best 0.433 (n = 1000) {[}0.402, 0.464{]}.

Proportions published in the companion archive whose per-cell replicate
counts are not reproduced in this archive (D6 0.633/0.733, D7
0.933/0.867, IC alternative-rule row 0.509/0.992) are quoted as
archived; their intervals are not recomputed here.

The misspecification cells sit one stage further along the diagnostic:
D6 (0.633/0.733) and D7 (0.933/0.867) are \emph{identifiable but
misspecified} --- the score finds a best module readily, and the price
is mechanism misattribution, the D2-axis failure shown in Table format
by the archived per-cell retention of the wrong module. Broader
candidate sets and their conditional operating characteristics remain an
open extension; nothing in the present archive bounds them.

\begin{center}\rule{0.5\linewidth}{0.5pt}\end{center}

\section{Supplement S2 --- the standard stated independently (two-page
specification and
checklist)}\label{supplement-s2-the-standard-stated-independently-two-page-specification-and-checklist}

\emph{Companion to paperF1\_retention\_framework\_v26\_restructured.md.
Page 1 is the normative specification; page 2 is the fillable checklist.
Nothing here depends on the worked examples of the main text.}

\subsection{Page 1 --- Normative
specification}\label{page-1-normative-specification}

\textbf{Scope.} A minimum reporting standard for a claim that a forecast
module is not retained: the claim must be produced by a pre-stated
decision rule, scoped by an information-set audit, and weighted by
measured operating characteristics.

\textbf{Obligations.}

\begin{enumerate}
\def\labelenumi{\arabic{enumi}.}
\tightlist
\item
  \textbf{Retention rule as an algorithm, fixed before scoring.} A
  forward-ordered ladder; declared baselines; a comparator map;
  horizons; a practical-equivalence margin with its basis stated
  (simulation-calibrated, decision-based, or pre-registered); outputs
  drawn from \{retained, not retained, declined on class grounds\}.
\item
  \textbf{Information-set audit.} One row per input quantity, classified
  available (dated at or before the origin), supplied (dated after it,
  provided regardless), or revised (dated before, published in a later
  vintage). The supplied-driver row determines whether the exercise is
  an operational forecast or a conditional hindcast.
\item
  \textbf{Operating-characteristic study, mandatory.} Synthetic series
  from known processes --- each in-class member as generating truth, a
  persistence-true null, and declared out-of-class processes --- with
  replicate counts, seeds, and adequacy thresholds fixed before any
  series is generated. Power, specificity, and mechanism misattribution
  are reported.
\end{enumerate}

\textbf{Output vocabulary.} retained / not retained / declined on class
grounds --- a substantive pre-gate, evaluated before scoring, declared
with reasons, not a threshold.

\textbf{Claim strength (certificate levels).} N0 --- the rule ran and
the verdict is its output. N1 --- plus archived margins and gate
decomposition. N2 --- plus the information-set audit. N3 --- plus
operating characteristics at the object's own length and noise attaining
the adequacy targets: the level at which non-retention may be reported
as evidence rather than description. A certificate expires on a revised
data vintage, a re-scored origin, or an amended rule; a combined claim
carries the lowest level of its parts.

\textbf{The worked rule (this article).} Origin-matched rolling-origin
RMSE at h = 1 and h = 5; H1 --- beat the next-simpler comparator by
strictly more than 5\%; H2 --- beat last-value persistence by strictly
more than 5\%; H3 --- both horizons; class-grounds pre-gate evaluated
first. The 5\% band is a practical-equivalence margin, fixed by
pre-registration.

\textbf{Reporting a verdict.} Margins; gate decomposition; descriptive
uncertainty (Diebold--Mariano and moving-block bootstrap ---
descriptive; verdicts do not rest on them); band and its basis;
certificate level; archives (pinned seeds, seed maps, every score and
forecast file).

\subsection{Page 2 --- Fillable
checklist}\label{page-2-fillable-checklist}

{\def\LTcaptype{none} % do not increment counter
\begin{longtable}[]{@{}
  >{\raggedright\arraybackslash}p{(\linewidth - 2\tabcolsep) * \real{0.5000}}
  >{\raggedright\arraybackslash}p{(\linewidth - 2\tabcolsep) * \real{0.5000}}@{}}
\toprule\noalign{}
\begin{minipage}[b]{\linewidth}\raggedright
Field
\end{minipage} & \begin{minipage}[b]{\linewidth}\raggedright
Entry
\end{minipage} \\
\midrule\noalign{}
\endhead
\bottomrule\noalign{}
\endlastfoot
Object: domain, predictand, period, origins (n) & \\
Target vintage status (revised / as-published) & \\
Ladder (forward-ordered) & \\
Baselines & \\
Comparator map & \\
Horizons & \\
Rule version; band; band basis & \\
Class-grounds declarations (with reasons) & \\
Information-set audit --- one row per quantity: available / supplied /
revised & \\
Margins versus baseline and comparator, per horizon & \\
Intervals (bootstrap / HAC) & \\
Gate decomposition per module (H1, H2, H3, pre-gate) & \\
Operating-characteristic study: DGPs, replicates, seeds, adequacy
thresholds & \\
Power / specificity / misattribution per cell & \\
Verdict per module & \\
Certificate level (N0--N3) & \\
Expiry conditions recorded & \\
Archive pointers (data, scores, seed maps) & \\
Sign-off: analyst, date & \\
\end{longtable}
}

The checklist is the required reporting form for every future
application registered in Section 8 of the main text. ---

\section{Supplement S3 --- fillable instruments: the information-set
audit canvas and the verdict
record}\label{supplement-s3-fillable-instruments-the-information-set-audit-canvas-and-the-verdict-record}

\emph{S3 is the fillable layer of Supplements S1--S2 of this supplement:
the audit table an analyst fills in before scoring, and the verdict
record the rule outputs at the end. Both are empty by design: they are
filled per application, never pre-filled.}

\subsection{S3.1 Information-set audit canvas (one row per driver
quantity)}\label{s3.1-information-set-audit-canvas-one-row-per-driver-quantity}

The columns of Section 3 of the main text, restated as a form. Classify
each quantity at each origin: \textbf{available} (dated at or before the
origin), \textbf{supplied} (dated after, provided regardless), or
\textbf{revised} (dated before, published only in a later vintage). The
supplied-driver row determines whether the exercise is an operational
forecast or a conditional hindcast.

{\def\LTcaptype{none} % do not increment counter
\begin{longtable}[]{@{}
  >{\raggedright\arraybackslash}p{(\linewidth - 12\tabcolsep) * \real{0.1429}}
  >{\raggedright\arraybackslash}p{(\linewidth - 12\tabcolsep) * \real{0.1429}}
  >{\raggedright\arraybackslash}p{(\linewidth - 12\tabcolsep) * \real{0.1429}}
  >{\raggedright\arraybackslash}p{(\linewidth - 12\tabcolsep) * \real{0.1429}}
  >{\raggedright\arraybackslash}p{(\linewidth - 12\tabcolsep) * \real{0.1429}}
  >{\raggedright\arraybackslash}p{(\linewidth - 12\tabcolsep) * \real{0.1429}}
  >{\raggedright\arraybackslash}p{(\linewidth - 12\tabcolsep) * \real{0.1429}}@{}}
\toprule\noalign{}
\begin{minipage}[b]{\linewidth}\raggedright
Quantity
\end{minipage} & \begin{minipage}[b]{\linewidth}\raggedright
Role
\end{minipage} & \begin{minipage}[b]{\linewidth}\raggedright
Dated at origin?
\end{minipage} & \begin{minipage}[b]{\linewidth}\raggedright
Published by origin?
\end{minipage} & \begin{minipage}[b]{\linewidth}\raggedright
Classification
\end{minipage} & \begin{minipage}[b]{\linewidth}\raggedright
Modules receiving it
\end{minipage} & \begin{minipage}[b]{\linewidth}\raggedright
Source / vintage note
\end{minipage} \\
\midrule\noalign{}
\endhead
\bottomrule\noalign{}
\endlastfoot
& & & & & & \\
& & & & & & \\
& & & & & & \\
& & & & & & \\
& & & & & & \\
& & & & & & \\
\end{longtable}
}

Reading of the filled table: - a \textbf{supplied} driver with
\texttt{role:\ driver} converts a forecast into a conditional hindcast
\emph{by record}, not by inference; - a \textbf{revised}
predictand/expiry row triggers the certificate's vintage-expiry
condition (Section 2.2 N-level expiry); - any module receiving a
post-origin quantity is excluded from the operational comparison and
reported separately.

\subsection{S3.2 Verdict record (the negative certificate,
human-fillable and
machine-readable)}\label{s3.2-verdict-record-the-negative-certificate-human-fillable-and-machine-readable}

Fill at the close of an application. The same record serialises to the
JSON schema of S3.3.

\begin{itemize}
\tightlist
\item
  \texttt{predictand} --- the quantity forecast (units).
\item
  \texttt{data\_vintage} --- the assessment/record vintage used.
\item
  \texttt{forecast\_origins} --- list of origins; count n.
\item
  \texttt{ladder} --- forward-ordered ladder rungs (frozen labels where
  applicable).
\item
  \texttt{rule\_version} --- the pre-registered rule identifier (R2 for
  the worked rule).
\item
  \texttt{band} and \texttt{band\_basis} --- the practical-equivalence
  margin and its basis (pre-registered / simulation-calibrated /
  decision-based, Section 8).
\item
  \texttt{benchmark} --- the baseline(s).
\item
  \texttt{comparators} --- the comparator map (next-simpler rung per
  module).
\item
  \texttt{horizons} --- the horizon set.
\item
  \texttt{loss} --- the scoring loss.
\item
  \texttt{information\_set} --- pointer to the filled S3.1 canvas.
\item
  \texttt{verdicts} --- one per module, drawn from \{retained, not
  retained, declined on class grounds\}; each with its margins, gate
  decomposition, and --- if declined --- the §3a class-grounds reason.
\item
  \texttt{operating\_characteristics} --- the class-conditional power
  under its own generating truth and the null specificity, with
  Monte-Carlo intervals (S1.5).
\item
  \texttt{certificate\_level} --- N0 / N1 / N2 / N3 (Section 2.2).
\item
  \texttt{expiry} --- the trigger conditions (new vintage, re-scored
  origin, amended rule); a combined claim carries the lowest level of
  its parts.
\item
  \texttt{archives} --- data, scores, scripts, seed maps, commit/DOI.
\item
  \texttt{analyst}, \texttt{date}, \texttt{sign-off}.
\end{itemize}

\subsection{S3.3 Machine-readable
schema}\label{s3.3-machine-readable-schema}

The record of S3.2 serialises to JSON (YAML-compatible). The companion
script \texttt{certificate\_schema\_S3.py} (this folder) validates a
filled record with the standard library only: required fields, the
frozen verdict vocabulary, the band-basis vocabulary, the N0--N3 levels,
and the rule that a declined module must carry its class-grounds reason.
A blank template is printed with \texttt{-\/-example}. The validator
checks the form of a claim, not its truth.

\begin{center}\rule{0.5\linewidth}{0.5pt}\end{center}

\section{Supplement S4 --- positioning against the forecast-comparison
and reporting
literature}\label{supplement-s4-positioning-against-the-forecast-comparison-and-reporting-literature}

\emph{What each neighbouring framework provides, and the three reporting
obligations this article adds on top. Descriptive, not evaluative: the
columns mark what a framework asks for, not whether it is ``better''.}

{\def\LTcaptype{none} % do not increment counter
\begin{longtable}[]{@{}
  >{\raggedright\arraybackslash}p{(\linewidth - 8\tabcolsep) * \real{0.2000}}
  >{\raggedright\arraybackslash}p{(\linewidth - 8\tabcolsep) * \real{0.2000}}
  >{\raggedright\arraybackslash}p{(\linewidth - 8\tabcolsep) * \real{0.2000}}
  >{\raggedright\arraybackslash}p{(\linewidth - 8\tabcolsep) * \real{0.2000}}
  >{\raggedright\arraybackslash}p{(\linewidth - 8\tabcolsep) * \real{0.2000}}@{}}
\toprule\noalign{}
\begin{minipage}[b]{\linewidth}\raggedright
Framework
\end{minipage} & \begin{minipage}[b]{\linewidth}\raggedright
What it tests / provides
\end{minipage} & \begin{minipage}[b]{\linewidth}\raggedright
Information-set audit
\end{minipage} & \begin{minipage}[b]{\linewidth}\raggedright
Equivalence margin
\end{minipage} & \begin{minipage}[b]{\linewidth}\raggedright
Class operating characteristics
\end{minipage} \\
\midrule\noalign{}
\endhead
\bottomrule\noalign{}
\endlastfoot
Diebold--Mariano (1995) & equal predictive accuracy, pair & -- & -- &
-- \\
White (2000), \emph{Econometrica} 68(5): 1097--1126 {[}1{]} & reality
check over a searched model universe (data-snooping control) & -- & -- &
-- \\
Hansen (2005), \emph{JBES} 23: 365--380 {[}2{]} & superior predictive
ability over a universe & -- & -- & -- \\
Hansen, Lunde \& Nason (2011), \emph{Econometrica} 79(2): 453--497
{[}3{]} & model confidence set --- the set containing the best with
given confidence & -- & -- & -- \\
Giacomini \& White (2006), \emph{Econometrica} 74(6): 1545--1578 {[}4{]}
& conditional predictive ability, rolling windows & -- & -- & -- \\
Clark \& West (2007), \emph{J. Econometrics} 138(1): 291--311 {[}5{]} &
equal accuracy for nested models & -- & -- & -- \\
Equivalence / non-inferiority testing, Wellek (2010), 2nd ed., Chapman
\& Hall/CRC {[}6{]} & equivalence within a pre-set margin & -- & margin
= inferential null & -- \\
Accuracy measures, Hyndman \& Koehler (2006), \emph{IJF} 22(4): 679--688
{[}7{]} & scaled error measures (MASE) & -- & -- & -- \\
\textbf{This article, R2} & (nothing new to estimate) & \textbf{typed
available/supplied/revised audit} & \textbf{practical-equivalence margin
on the decision} & \textbf{class-conditional power + null
specificity} \\
\end{longtable}
}

\subsection{Reading the table}\label{reading-the-table}

The right-hand columns are the three obligations of Section 2. The
multiple- comparison literature (White, Hansen, Hansen--Lunde--Nason)
corrects the \emph{selection} problem the rule runs into when many
modules are scored --- the reality check and the model confidence set
quantify how much of a ``winner'' is search. The nested-model test
(Clark--West) is exactly the comparator-gate problem stated as inference
rather than decision. The equivalence-testing literature (Wellek)
supplies the inferential form of the practical-equivalence margin this
article uses as a decision margin. None of them asks what a forecast was
allowed to see (the audit), treats the margin as an input to the
decision rule rather than the null itself, or asks whether the procedure
could have detected the rejected class (the operating characteristics).
Those three gaps are the article's only claims of novelty; the rest is
assembly.

\textbf{Relationship, not replacement.} The rule's multiple-testing
behaviour is \emph{descriptive here} --- the comparison is over a
declared ladder, not a searched universe --- but a reader porting the
standard to a wide model search should pair it with a reality-check or
model-confidence-set correction. That pairing is a stated route to
endorsement, not a defect of neighbouring work.

\subsection{Sources}\label{sources}

\begin{enumerate}
\def\labelenumi{\arabic{enumi}.}
\tightlist
\item
  White, H. (2000). A reality check for data snooping.
  \emph{Econometrica}, 68(5), 1097--1126.
\item
  Hansen, P. R. (2005). A test for superior predictive ability.
  \emph{Journal of Business \& Economic Statistics}, 23, 365--380.
  doi:10.1198/073500105000000063
\item
  Hansen, P. R., Lunde, A., \& Nason, J. M. (2011). The model confidence
  set. \emph{Econometrica}, 79(2), 453--497. doi:10.3982/ECTA5771
\item
  Giacomini, R., \& White, H. (2006). Tests of conditional predictive
  ability. \emph{Econometrica}, 74(6), 1545--1578.
  doi:10.1111/j.1468-0262.2006.00718.x
\item
  Clark, T. E., \& West, K. D. (2007). Approximately normal tests for
  equal predictive accuracy in nested models. \emph{Journal of
  Econometrics}, 138(1), 291--311.
\item
  Wellek, S. (2010). \emph{Testing statistical hypotheses of equivalence
  and noninferiority} (2nd ed.). Chapman \& Hall/CRC. ISBN
  978-1439808184.
\item
  Hyndman, R. J., \& Koehler, A. B. (2006). Another look at measures of
  forecast accuracy. \emph{International Journal of Forecasting}, 22(4),
  679--688. doi:10.1016/j.ijforecast.2006.03.001
\end{enumerate}

\begin{center}\rule{0.5\linewidth}{0.5pt}\end{center}

\subsection{S1.4 Reproducibility
package}\label{s1.4-reproducibility-package}

The working-example margins, the operating-characteristic cells (power,
specificity, mechanism-misattribution rows), the likelihood-ratio
derivations, and the calibrated replacement band recompute from the
registered campaigns below, archived at
https://github.com/MIKEAA2020/general-sustainability; checksums (md5,
first 8 hex) pin the release state. The pre-registered 10,000-replicate
archives are the dataset of record: alternative decision rules (Section
6.5) and all sensitivity variants (S1.1) are evaluated on the same
archived replicates without refitting, so the reported cells are
deterministic functions of registered files. Re-executing the campaigns
in a fresh environment regenerates every archived result file byte for
byte.

{\def\LTcaptype{none} % do not increment counter
\begin{longtable}[]{@{}
  >{\raggedright\arraybackslash}p{(\linewidth - 6\tabcolsep) * \real{0.2500}}
  >{\raggedright\arraybackslash}p{(\linewidth - 6\tabcolsep) * \real{0.2500}}
  >{\raggedright\arraybackslash}p{(\linewidth - 6\tabcolsep) * \real{0.2500}}
  >{\raggedright\arraybackslash}p{(\linewidth - 6\tabcolsep) * \real{0.2500}}@{}}
\toprule\noalign{}
\begin{minipage}[b]{\linewidth}\raggedright
Script (repo path)
\end{minipage} & \begin{minipage}[b]{\linewidth}\raggedright
Role in the paper
\end{minipage} & \begin{minipage}[b]{\linewidth}\raggedright
md5
\end{minipage} & \begin{minipage}[b]{\linewidth}\raggedright
\end{minipage} \\
\midrule\noalign{}
\endhead
\bottomrule\noalign{}
\endlastfoot
\texttt{framework/campaign\_cod\_t71\_band\_calibration.py} & Primary
simulation-calibrated band campaign on the cod T = 71 object (registered
prospective replacement band) & \texttt{b261aef5} & calibration
outputs \\
\texttt{framework/campaign\_cod\_t71\_band\_calibration\_checkpointed\_20260917.py}
& Checkpointed re-run of the same campaign (byte-identical checkpoints)
& \texttt{d1336ea2} & checkpoint files \\
\texttt{framework/campaign\_edwards\_band\_calibration.py} &
Simulation-calibrated band campaign on the Edwards J-17 object &
\texttt{26bb9780} & calibration outputs \\
\texttt{framework/campaign\_edwards\_band\_calibration\_audit.py} &
Independent audit of the Edwards calibration (recomputes rows) &
\texttt{b261a0ca} & audit output \\
\texttt{framework/certificate\_schema\_S3.py} & S3 certificate schema
validator for the audit layer & \texttt{19a22c18} & schema/validation
records \\
\texttt{phase\_c/campaign\_power.py} & Core power campaign (D1--D5
replicates) & \texttt{352d4e84} & replicate archives \\
\texttt{phase\_c/campaign\_misspecified.py} & Amendment-1 campaign
(D6--D7 out-of-class truth) & \texttt{ea603471} & replicate archives \\
\texttt{phase\_c/campaign\_origins.py} & Origin-set controls (mixed
versus origin-matched) & \texttt{0d3b2cf0} & origin tables \\
\texttt{phase\_c/campaign\_snr.py} & Signal-to-noise grid campaign &
\texttt{2c4a496e} & SNR tables \\
\texttt{phase\_c/campaign\_smoother.py} & Smoothed-series control
campaign & \texttt{ac2331ea} & smoother tables \\
\texttt{phase\_c/o6\_cod\_band\_calibration\_20260913.py} & Earlier cod
band-calibration layer (registered 2026-09-13) & \texttt{a5e21be2} & o6
outputs \\
\texttt{phase\_c/o9\_ic\_coprimary\_20260913.py} & IC co-primary
comparison layer & \texttt{e327f256} & o9 outputs \\
\texttt{phase\_c/analysis\_d67\_gain.py} & D6/D7 retention-gain analysis
(mean gains +3.1/+9.2 kt) & \texttt{fa367856} & gain tables \\
\texttt{phase\_c/analysis\_gate\_mcs.py} & Gate-attribution and
model-confidence-set analyses & \texttt{6cbdcf7c} & gate/MCS tables \\
\texttt{phase\_c/analysis\_regime\_ident.py} & Regime-identification
analysis & \texttt{a5186de2} & identification tables \\
\texttt{batch\ 7\ (audits\ of\ agent\ arena\ 1\ paper\ rewrites)/campaign\_e1\_dm\_uncertainty.py}
& Cod DM/bootstrap post-freeze layer (Spec A/B) & \texttt{81c559b8} &
results/e1\_dm\_uncertainty.csv \\
\texttt{batch\ 7\ (audits\ of\ agent\ arena\ 1\ paper\ rewrites)/campaign\_e1\_dm\_uncertainty\_regime.py}
& Cod regime-stratified DM layer & \texttt{6ced3555} &
results/e1\_dm\_uncertainty\_regime.csv \\
\texttt{batch\ 7\ (audits\ of\ agent\ arena\ 1\ paper\ rewrites)/campaign\_e3\_dm\_uncertainty.py}
& Edwards DM/bootstrap independent replication & \texttt{30963493} &
results/e3\_dm\_uncertainty.csv \\
\end{longtable}
}

\emph{Design-time exclusions.} Candidates rejected before scoring ---
including the naive regime-switch generator (900 → 888 → \ldots{} → 0 in
four steps, rejected at design time) --- are excluded from the
registered package by construction and are recorded in Section 6.2 of
the carrier.

\end{document}
